\section{Introduction}\label{sec:introduction}

This document revises the exposition of a proposed proof in the supplied
manuscript. George Coote's purpose in working with GPT 6.1 Sol was to make
that argument readable enough for him and others to examine its
mathematical validity without having to reconstruct numerous omitted
steps. The main statements and proof strategy were already present in
the supplied text. The revision expands definitions, calculations,
references, and logical transitions; it does not claim that Coote
originated the result or its proof, or that the proposed proof has been
independently verified. The following introduction describes the claim
and the argument presented for it.

The hot-spots problem arose from Rauch's discussion of the long-time
behavior of heat flow with insulating boundary conditions \cite{Rau75}.
Its spectral formulation asks whether the extrema of a first-positive
Neumann eigenfunction occur on the boundary. Ba\~nuelos and Burdzy
distinguished strict and nonstrict formulations, and assertions about
every eigenfunction from the existence of one favorable eigenfunction
\cite{BanuelosBurdzy}. These distinctions matter when the eigenvalue is multiple.

Let $\Omega\subset\mathbb R^2$ be a nonempty bounded simply connected open
set with $C^\infty$ boundary. In particular, $\Omega$ is connected. Set
\begin{equation}\label{eq:first-neumann-eigenvalue}
 \lambda_N=\inf_{\substack{0\ne v\in H^1(\Omega;\mathbb R)\\
                            \int_\Omega v=0}}
 \frac{\int_\Omega|\nabla v|^2}{\int_\Omega v^2},
\end{equation}
and let $V$ be the real Neumann eigenspace for $\lambda_N$.
This is the first positive eigenvalue; sources that count the constant
eigenfunction first call it the second Neumann eigenvalue.

The supplied manuscript claims the following theorem.
\begin{samepage}
\begin{theorem}\label[theorem]{thm:strict-hot-spots}
For every $0\ne u\in V$,
\[
 \nabla u(x)\ne0\qquad(x\in\Omega).
\]
In particular,
\begin{equation}\label{eq:strict-hot-spots}
 \min_{y\in\partial\Omega}u(y)<u(x)<\max_{y\in\partial\Omega}u(y)
 \qquad(x\in\Omega).
\end{equation}
\end{theorem}
\end{samepage}

The theorem imposes no convexity or symmetry condition. It applies to
every member of the first positive eigenspace, including when the
eigenvalue is multiple. The nonvanishing-gradient conclusion also rules
out interior saddle points.

\paragraph{Background and methods.}
Topology matters in the hot-spots problem. Burdzy and Werner constructed
a planar counterexample with two holes and a simple first-positive
eigenvalue \cite{BW99}. Bass and Burdzy obtained examples with both extrema
strictly in the interior \cite{BB00}, and Burdzy subsequently obtained
this behavior with just one hole \cite{Bur05}. The simply connected
planar conjecture is stated explicitly in \cite[Conjecture 1.2(ii)]{Bur05}.

Several geometric special cases already establish strong gradient
conclusions. Jerison and Nadirashvili proved a positive directional
derivative for every first-positive eigenfunction on convex planar domains
with two orthogonal reflection symmetries, allowing multiple eigenvalues
\cite[Theorem 1.4]{JN00}. Probabilistic coupling methods developed by
Ba\~nuelos and Burdzy \cite{BanuelosBurdzy} led to further results. Pascu
proved boundary-only extrema for antisymmetric eigenfunctions on
$C^{1,\alpha}$ convex planar domains with one reflection symmetry, where
$0<\alpha<1$ \cite{Pas02}. Atar and Burdzy treated every first-positive
eigenfunction on lip domains, a class of Lipschitz domains bounded between
graphs of functions with Lipschitz constants at most one \cite{AB04}.
Miyamoto obtained interior critical-point exclusion under a
spectral-diameter condition, giving nonsymmetric nearly circular convex
examples \cite[Lemma 1.2 and Theorem A]{Miy07}.

For polygons, Siudeja proved hot-spots results for a class of acute
triangles \cite{Siu15}. Judge and Mondal, with their subsequent erratum,
excluded interior critical points for every first-positive Neumann
eigenfunction on every Euclidean triangle \cite{JM20,JM22}. Chen, Gui and
Yao established the finer classification of nonvertex critical points
\cite{CGY26}. These results concern domains with corners;
\cref{thm:strict-hot-spots} concerns smooth simply connected planar domains.

Recent quantitative work measures how far the hot-spots conclusion can
fail. For general bounded connected Lipschitz domains, de Dios Pont, Hsu
and Taylor determined the sharp dimension-dependent upper bound for the
ratio of a first-positive Neumann eigenfunction's supremum over the domain
to its maximum on the boundary \cite[Definition 1 and Theorem 4]{dDPHT25}.
Deng, Jiang and Yang obtained quantitative bounds for convex domains in
two-dimensional space forms \cite{DJY26}. Such ratios concern extrema;
excluding every interior critical point also requires excluding saddle points.

The supplied proof starts from Rohleder's variational principle for tangent vector
fields \cite{Rohleder,Roh24}. In its lip-domain application, taking
componentwise absolute values after a suitable rotation preserves the
required tangent boundary condition and yields directional monotonicity.
Its argument uses scalar boundary multipliers to keep tangency on an
arbitrary smooth boundary, constructing those multipliers through a
reciprocal Green-kernel theorem.

Each Cartesian derivative of an eigenfunction solves the same Helmholtz
equation, but its boundary values need not have a fixed sign. The Neumann
condition constrains the two derivatives together: their vector is tangent
to the boundary. The argument uses that constraint through a nonnegative quadratic
form on tangent boundary data. Its nullspace consists exactly of the
boundary traces of first-positive eigenfunction gradients. The kernel
construction supplies scalar boundary multipliers that preserve this nullspace if an
interior critical point exists.

\paragraph{The mechanism.}
Write $\mathbb D=\{z\in\mathbb C:|z|<1\}$ and $S^1=\partial\mathbb D$,
and let $\Phi:\overline{\mathbb D}\to\overline\Omega$ be a smooth
conformal parametrization. For a tangent vector field $X$, the
divergence--curl energy $Q(X)$ is at least $\lambda_N\|X\|_2^2$, with equality
exactly for gradients of functions in $V$; see \cref{prop:vector-variational}.
After the conformal change of variables, solving the scalar Helmholtz
equation in each fixed Cartesian component extends Lipschitz boundary
data $h:S^1\to\mathbb R^2$ to a vector field $X_h$ on $\Omega$.
For tangent data, the quadratic form
\[
 E(h)=Q(X_h)-\lambda_N\|X_h\|_2^2
\]
is nonnegative, and it vanishes exactly when $X_h=\nabla v$ for some
$v\in V$. Here tangency at $s$ means tangency to $\partial\Omega$ at
$\Phi(s)$. For a nonzero eigenfunction $u\in V$, put
$g(s)=\nabla u(\Phi(s))$. This is the gradient in the coordinates of
$\Omega$, parametrized on the circle. Its extension is $X_g=\nabla u$,
so $E(g)=0$, and a real Lipschitz scalar multiplier $b$ satisfies
\[
 E(bg)=\frac12\iint_{S^1\times S^1}
 N(s,t)(b(s)-b(t))^2g(s)\cdot g(t)\,ds\,dt.
\]
Here $N$ is a symmetric boundary kernel, strictly positive off the
diagonal, and the corresponding interior
evaluation kernel is denoted by $K$; see \cref{prop:multiplier-identity}.

The sign of $g(s)\cdot g(t)$ prevents a direct positivity argument in
this identity. The kernel theorem supplies multipliers whose squared
differences remove the factor $N$. For every interior point $p$,
\[
 D_p(s,t)=\frac{K(p,s)K(p,t)}{N(s,t)}\quad(s\ne t),
 \qquad D_p(s,s)=0,
\]
is conditionally negative semidefinite. \Cref{lem:hilbert-embedding}
constructs a separable real Hilbert space and an injective Lipschitz map
$b$ from the circle into that space, with squared distances $D_p$.
The construction uses the positive anchored kernel, a quotient by the
zero-norm subspace, and completion; the kernel estimates give injectivity
and the Lipschitz bound. Applying the energy identity to the coordinates
$b_j$ of $b$ in an orthonormal basis, and justifying the sum by dominated
convergence, gives
\[
 \sum_j E(b_jg)=\frac12|\nabla u(\Phi(p))|^2.
\]
At a critical point every nonnegative summand vanishes. The equality case
then gives $v_j\in V$ with boundary gradient $b_jg$. If the set
$S_g=\{s:g(s)\ne0\}$ is dense in the circle, injectivity of $b$ produces
three linearly independent eigenfunctions, contradicting
\cref{prop:multiplicity}. If $S_g$ is not dense, the gradients vanish
on an open boundary arc. Two independent eigenfunctions then have a
nonzero linear combination with both zero value and zero normal
derivative on that arc. Its local extension by zero solves the Helmholtz
equation across the arc, and interior uniqueness gives the contradiction.

The construction of $D_p$ is the central analytic step. It first
regularizes the disk Green function so that the entrywise reciprocal of
every finite kernel matrix has at most one positive eigenvalue. This
property survives successive rank-one updates along columns of the
current kernel matrix, which sum the Green resolvent over a finite set of
integration nodes. Positive quadrature and monotone limits then pass from
those matrices to the singular Green kernel and finally to its boundary
kernels. Taking the boundary limit before removing a compact-support
cutoff avoids treating a Poisson kernel as an $L^2$ function when it need
not be one.

The kernel theorem, \cref{thm:squared-distance}, holds for every bounded
nonnegative disk potential satisfying a weighted Dirichlet inequality
with constant $d<1$, independently of the conformal origin of that
potential. The conformal potential satisfies this condition by
\cref{eq:weighted-dirichlet}. The reciprocal-matrix condition is related
to the McCullough--Quiggin characterization of complete Nevanlinna--Pick
kernels \cite{McC92,McC94,Qui93,Qui94}; see also
\cite[Corollary 1.12]{AglerMcCarthy}. The condition used here requires at most one positive
eigenvalue in each finite sampled reciprocal matrix, allowing zero
diagonal entries in the limiting matrix. This condition concerns finite
matrices, rather than reciprocals of integral operators.
Schoenberg's squared-distance correspondence underlies the Hilbert-space
construction proved in \cref{lem:hilbert-embedding} \cite{Sch38}.
The analytic work here concerns the singular Green kernel, its
positive-potential resolvent, and its boundary limit.

\paragraph{Organization and background.}
\Cref{sec:preliminaries} proves the spectral gap, the smooth conformal
reduction, and the multiplicity bound. \Cref{sec:vector-fields} proves
the vector variational principle. \Cref{sec:boundary-kernels} constructs
the boundary kernels and derives the multiplier identity.
\Cref{sec:reciprocal-green} proves their conditional negativity, and
\Cref{sec:hilbert-multipliers} completes the argument.
The argument relies on Sobolev compactness, Poincar\'e and trace inequalities on smooth
bounded domains, Poisson solvability and elliptic boundary regularity,
the maximum and boundary point principles, elementary complex analysis,
and planar separation. The later sections give detailed arguments for
the specialized ingredients.
The smooth-boundary regularity input is
\cite[Chapter III, Theorems 6.1, 6.4 and 6.5]{Showalter}: $L^2$ Poisson
data give $H^2$ solutions under homogeneous Dirichlet or Neumann
conditions, and higher regularity follows with smoother data.
Subtracting a smooth extension handles the inhomogeneous Dirichlet data
used below. Bootstrap and Sobolev embedding make weak Neumann
eigenfunctions smooth up to the boundary. Interior solutions of
$(\Delta+\lambda_N)v=0$ are analytic; alternatively,
$e^{\sqrt{\lambda_N}t}v(x)$ is harmonic in one additional variable,
as shown in \cref{fact:helmholtz-uniqueness}.

\paragraph{Human revision and use of GPT 6.1 Sol.}
The exposition in \cref{sec:preliminaries,sec:vector-fields,sec:boundary-kernels,sec:reciprocal-green,sec:hilbert-multipliers}
was developed through a passage-by-passage revision directed by George
Coote, using OpenAI's GPT 6.1 Sol. The aim was to reduce the effort
needed to understand and review the supplied proof. Coote
selected the scope and pace of the work, identified compressed or
ambiguous arguments, challenged terminology and unsupported transitions,
and decided which explanations and calculations should enter the text.
The revision consistently uses $\lambda_N$ for the first positive
Neumann eigenvalue and makes the test spaces, trace conditions, equality
cases, and dependencies between results explicit.

In \cref{sec:preliminaries}, this scrutiny led to fuller explanations of
variational equality and plane-wave independence, smooth conformal
boundary regularity, the boundary-sign argument, and the multiplicity
bound. Coote required the planar separation argument to rest
on a precise crosscut statement and proof, with the figure serving as an
illustration. In \cref{sec:vector-fields}, the revisions explain the
construction and compatibility of the Neumann and Dirichlet potentials,
the truncation and weak-gradient arguments, Sobolev integration by parts,
and the equality case. They also expand the curvature identity and its
passage from smooth fields to Sobolev data.

In \cref{sec:boundary-kernels}, Coote requested explicit
conformal changes of variables, weak solution and uniqueness arguments,
kernel bounds, and justifications for exchanging integrals and sums.
The revised energy calculation identifies the componentwise extensions,
the Fourier normalization, curvature cancellation, and absolute
integrability in the multiplier identity. In \cref{sec:reciprocal-green},
Coote required a clear distinction between finite sampled
matrices and integral operators, a literature comparison for the
reciprocal condition, and explicit quadrature nodes, weights, matrix
updates, and convergence arguments. The revision separates diagonal
from off-diagonal boundary limits and explains the compact-support
cutoff before its removal. In \cref{sec:hilbert-multipliers}, Coote
requested the Hilbert-space construction, pointwise boundary
gradient identities, the transfer of linear independence to
eigenfunctions, and the integration-by-parts calculation for local zero
extension. The final extrema argument explicitly uses the extreme value
theorem and Fermat's stationary-point theorem from elementary calculus:
at an interior local maximum or minimum, a differentiable function has
zero gradient.

GPT 6.1 Sol supplied much of the new mathematical prose, detailed
derivations, supporting calculations, reference work, and LaTeX
implementation in response to these interventions. Coote
reviewed the proposed explanations, controlled their level of detail,
and authorized successive additions. This account summarizes the
section-by-section contribution records accompanying the source files;
it describes the development of the revised exposition. The author line
identifies Coote's work on this revision with GPT 6.1 Sol; it is not a
claim that he discovered or proved the underlying result. The recorded
mathematical checks were model-assisted reviews of the arguments. The
purpose of making those arguments explicit is to enable further
mathematical review, and their correctness remains to be assessed.
